\chapter{Literature Review}
\label{ch:literature_review}

\section{LLM Token Economics}
\label{sec:token_economics}

The cost structure of large language model APIs is fundamentally different from the cost
structure of traditional database queries or REST services. Relational database calls are
priced by compute time, which correlates weakly with the size of the query or its result.
LLM API calls are priced by the token, which creates a direct linear relationship between
the amount of context supplied and the monetary cost of the call. Brown et al.~\citep{brown2020gpt3}
established the scaling properties that make large models capable, and those same scaling
properties---more parameters, larger context windows---are what make the token-per-call
model so costly at enterprise scale.

The context window problem is not new. Early transformer models were limited to 512 or
1,024 tokens; modern models accept 128,000 or more. Paradoxically, larger context windows
have made the cost problem worse, not better, because they enable---and in production
settings, encourage---practitioners to dump entire datasets into a single prompt. Vaswani
et al.~\citep{vaswani2017attention} showed that the self-attention mechanism scales
quadratically with sequence length, and while linearized attention variants have reduced
the computational cost, the billing models of commercial APIs remain roughly linear in the
number of tokens processed. The result is that a team that migrates from a 4K-context
deployment to a 128K-context deployment does not necessarily save money; it often spends
more because each query now ingests far more data than is strictly necessary.

Retrieval-Augmented Generation~\citep{lewis2020rag} addresses one version of this problem
by indexing documents into a vector store and retrieving only the top-$k$ most relevant
chunks at query time. RAG has been widely adopted and is effective for question-answering
over static document corpora. Its limitation for the present problem is that supply chain
event streams are not static corpora. Events arrive continuously, their relevance is
determined by recency and by structural position within the business graph, not by semantic
similarity to a query string, and the quantity of interest is often what \emph{changed},
not what \emph{exists}. A vector similarity search between the query ``what are the
critical supply risks today?'' and a raw event record will not reliably surface the fact
that a supplier's centrality in the graph increased because a new contract made it the
sole source for three products.

Naive context chunking---splitting the event stream into fixed-size windows and forwarding
each window independently---avoids the RAG dependency but discards relationship information.
A chunk of 100 sales events contains no indication that two of the suppliers involved are
linked to the same warehouse, or that an anomaly in one product's defect rate correlates
with a delay pattern in a geographically distant warehouse. The present project treats
this structural relationship information as the primary signal and builds the compression
strategy around preserving it.

\section{Knowledge Graphs for Business Intelligence}
\label{sec:knowledge_graphs}

Knowledge graphs have a long history in enterprise data management, predating the current
LLM wave by at least two decades. The Resource Description Framework (RDF) and the Web
Ontology Language (OWL) established formal foundations for representing entities,
relationships, and inference rules over structured business data~\citep{berners2001semantic}.
Industrial deployments at companies like Google, Amazon, and LinkedIn demonstrated that
graph-based representations of product catalogs, supply networks, and social connections
support query patterns that relational schemas handle poorly, particularly multi-hop
relationship queries and anomaly detection via structural deviation.

For supply chain modeling, graph representations have proven especially natural. A supplier
node connected by ``supplies'' edges to product nodes, which are in turn connected by
``stored\_at'' edges to warehouse nodes and by ``sold\_in'' edges to region nodes, captures
the essential topology of the business in a form that is both human-readable and
algorithmically tractable. Centrality measures on this graph have direct business
interpretations. A supplier with high betweenness centrality sits on many shortest paths
between other nodes; a delay from that supplier disrupts more downstream processes than
a delay from a peripheral supplier. Degree centrality measures how many direct connections
a node has; a warehouse with high degree is a hub whose stockout affects more product lines
simultaneously.

NetworkX~\citep{hagberg2008networkx} provides the practical implementation substrate for
this approach. Its Python API supports the full range of graph operations needed for this
project: node and edge creation with arbitrary attribute dictionaries, snapshot export for
delta computation, and the centrality algorithms described above. The library's memory
efficiency for graphs of the scale used here---30 nodes and roughly 350--400 edges---is
more than adequate; the entire graph fits within a few kilobytes of RAM and is reconstructed
from scratch on each event batch in milliseconds.

Graph-based anomaly detection is a well-studied area. Methods range from community
detection via the Louvain algorithm~\citep{blondel2008louvain} to embedding-based
approaches that map nodes to vector spaces and flag statistical outliers. The present
project uses a simpler but interpretable approach: an \texttt{anomaly\_score} attribute
on each node is updated by the event generator when anomalous conditions are detected
(quality defect rates above a threshold, delay severity above a threshold), and the delta
engine surfaces changes in this attribute as high-saliency events. This keeps the
anomaly detection transparent and audit-friendly, which matters in enterprise supply
chain contexts where decisions based on automated flags must be explainable.

\section{Saliency Scoring and Information Filtering}
\label{sec:saliency}

Saliency, in the sense of measuring which parts of an input are most relevant for a
downstream decision, has been a core concept in information retrieval since the early work
on term frequency--inverse document frequency (TF-IDF) weighting~\citep{salton1988tfidf}.
The intuition is the same in both contexts: not all elements of an input contribute equally
to the output, and a good scoring function can identify the high-contribution elements
without evaluating the output for each possible subset.

Attention mechanisms in neural networks~\citep{bahdanau2015attention} operationalized this
idea within deep learning, allowing a model to learn which parts of the input to ``attend
to'' for a given prediction. The additive attention formulation of Bahdanau et al.\ and the
scaled dot-product attention of Vaswani et al.~\citep{vaswani2017attention} are the
architectural foundations of modern transformer models. This project does not train an
attention model; instead, it uses a hand-crafted saliency function that combines three
structural signals (degree centrality, betweenness centrality, and attribute magnitude)
into a scalar score computable in microseconds per node without any gradient computation.

Threshold-based filtering on saliency scores is well-established in information retrieval.
Buckley and Voorhees~\citep{buckley2000precision} showed that aggressive top-$k$ retrieval
with appropriate thresholds preserves most of the recall of a full-corpus search while
cutting the number of documents processed by an order of magnitude. This project applies
the analogous result to structured events: score each change by its business impact,
discard changes below a saliency threshold, and forward only the high-impact changes to
the agent. Empirical tuning on the ground truth anomaly patterns found that a threshold of
$\theta = 0.4$ on the composite saliency score preserved 93.8\% recall while filtering
roughly 90--95\% of routine low-importance events.

The challenge in defining saliency for supply chain events is that importance is contextual.
A price change of 5\% is routine noise in a volatile commodity market but a significant
signal in a contract-bound manufacturing context. The present project handles this
by incorporating the magnitude of attribute changes as the third term in the saliency
formula, giving larger attribute changes higher scores regardless of baseline volatility.
This is a deliberate simplification; a production system would benefit from domain-specific
volatility normalization, which is identified as a future work item in
Chapter~\ref{ch:conclusion}.

\section{Multi-Agent LLM Systems}
\label{sec:multi_agent}

Multi-agent architectures for LLM-based systems have attracted substantial attention since
the release of frameworks like LangChain~\citep{chase2022langchain} and
CrewAI~\citep{crewai2024}. The canonical pattern involves assigning specialized roles to
individual agents---an analyst agent, a risk assessment agent, a forecasting agent---and
routing outputs between them through an orchestrator. Each agent receives a context
specific to its role and returns a structured response that other agents can use.

The economics of multi-agent systems are worse than single-agent systems by a multiplier
equal to the number of agent invocations. If a pipeline calls three agents per user
query, token costs are tripled. If each agent receives the full raw context rather than
a compressed briefing, costs are tripled again. This multiplicative effect is the primary
driver of the compression requirement: in a three-agent pipeline, achieving 10$\times$
compression at the input stage reduces total cost by an order of magnitude. The 87$\times$
compression measured in this project at 1,000 events reduces three-agent costs by nearly
two orders of magnitude relative to the naive full-context approach.

Tool-use and function-calling APIs have partly addressed the retrieval problem within
agent architectures. An agent that can call a database query tool or a structured API
does not need to receive the entire database in its context; it can fetch only what it
needs at inference time. The limitation is that tool-use adds latency---each tool call is
a round-trip to an external service---and requires the agent to know in advance what to
fetch, which is not always possible for exploratory analytical queries over live event
streams.

The present project's approach is complementary to tool-use. The distillation pipeline
runs as a preprocessing step before any agent is invoked, producing a briefing that
captures the most salient state of the business at the current moment. An agent receiving
this briefing can still use tools for follow-up queries, but the primary context it needs
to answer most operational questions is already present in the briefing without additional
round-trips.

\section{Gap Analysis}
\label{sec:gap_analysis}

A review of the approaches above reveals a consistent pattern: each existing method
addresses one dimension of the context compression problem but not all three simultaneously.
RAG addresses the retrieval dimension but not the structural dimension (relationship inference)
or the change dimension (what changed since the last snapshot). Summarization approaches
condense long texts but operate on unstructured natural language and do not exploit the
graph structure of business data. Embedding-based filtering addresses semantic similarity
but, like RAG, does not naturally capture structural change in a dynamic graph.
Table~\ref{tab:gap_analysis} summarizes this comparison.

\begin{table}[H]
\centering
\caption{Comparison of context compression approaches for LLM pipelines}
\label{tab:gap_analysis}
\begin{tabular}{lcccc}
\toprule
\textbf{Method} & \textbf{Token Reduction} & \textbf{Quality Maintained} & \textbf{Structured Output} & \textbf{Open Source} \\
\midrule
RAG (BM25 / embedding) & High & Moderate & No & Yes \\
LLM summarization & Moderate & Variable & No & Yes \\
Embedding-based filtering & Moderate & Moderate & No & Yes \\
\textbf{This work} & \textbf{87$\times$ max} & \textbf{93.8\% recall} & \textbf{Yes} & \textbf{Yes} \\
\bottomrule
\end{tabular}
\end{table}

The gap this project fills is the absence of a system that combines all three: symbolic
delta computation over a live knowledge graph, neural saliency scoring of those deltas,
and structured distillation into an LLM-ready briefing with explicit severity
classifications. No existing open-source package or published paper describes this
combination in an end-to-end benchmarked implementation. The present work is, to the
author's knowledge, the first open system to report token compression and recall metrics
for this architecture on a supply chain domain.
